\begin{abstracteng}
\tl{Automatic placenta segmentation in magnetic resonance imaging (MRI) is a fundamental step for
quantitative pregnancy analysis and for the systematic use of medical imaging\-
information. The task is particularly challenging due to strong variability in
organ\- morphology and appearance, ambiguous boundaries, and the limited
availability of sufficiently annotated data.}

\tl{This thesis performs a systematic comparative evaluation of modern deep
learning architectures\- for 3D placenta segmentation. A unified,
reproducible experimental \-framework was developed in MONAI/PyTorch,
integrating preprocessing, sampling, \-augmentation, training, and evaluation, so
that the comparison is fair and well documented.}

\tl{The study covers U-Net-type architectures, Transformer-based
approaches, and state-space-based models, assessed through Dice/IoU
and qualitative inspection of the \-produced masks. The findings highlight that
strict pipeline standardization is a critical factor for reliable
conclusions, and that different architecture families exhibit distinct strengths
in stability, boundary quality, and overall segmentation behavior.}

\tl{Overall, the thesis provides a clear reference framework for choosing
placenta \-segmentation models in MRI. At the same time, it underlines the
need for further investigation of generalization on independent datasets,
external validation, and the integration of uncertainty estimation techniques.}

   \begin{keywordseng}
    \tl{MRI, Placenta, Medical image segmentation, Deep Learning, MONAI,
    Comparative evaluation\-, U-Net, Transformers, State Space Models,
    Reproducibility}
   \end{keywordseng}

\end{abstracteng}
